% ============================================================
% chapter2-Search_Pipeline.tex
% ============================================================
\chapter{La Pipeline di Elaborazione della Query}

La Search Pipeline costituisce il nucleo operativo dell'assistente ed è
orchestrata interamente all'interno del modulo \path{search_pipeline.py}.
Alla ricezione di una richiesta in linguaggio naturale, il sistema attiva
una sequenza strutturata di elaborazioni che sfrutta la concorrenza asincrona
offerta da \path{asyncio}, combinando interrogazioni parallele a database
interni ed API esterne per produrre risultati classificati, arricchiti e
contestualizzati. La pipeline si articola in dodici fasi distinte, ciascuna
delle quali alimenta le successive con strutture dati progressivamente più
raffinate.

% ------------------------------------------------------------------
\subsection*{Fase 1: Recupero Anticipato della Memoria (Phantom Fetch)}
% ------------------------------------------------------------------
Non appena la richiesta raggiunge la pipeline, viene avviata in background
un'interrogazione su \textbf{Qdrant}, il database vettoriale del sistema,
tramite il metodo \path{retrieve_context}. Questa operazione, denominata
\emph{Phantom Fetch}, non attende le fasi di analisi semantica: agisce sul
testo grezzo per recuperare, con il massimo anticipo temporale, le inserzioni
già validate in sessioni precedenti su query semanticamente affini. Il
principio è semplice: prodotti già incontrati e giudicati positivamente
costituiscono una baseline di alta qualità, e il loro recupero anticipato
parallelizza la latenza con le fasi computazionalmente più intensive.

% ------------------------------------------------------------------
\subsection*{Fase 2: Recupero dello Storico Conversazionale}
% ------------------------------------------------------------------
Per mantenere coerenza su conversazioni multi-turno, la pipeline recupera
la cronologia recente della sessione tramite \path{redis_client.get_user_queries},
estraendo le ultime tre query (\path{history[:3]}). Questo \textit{snapshot}
contestuale è volutamente leggero: fornisce un riferimento minimo per
risolvere anafore e riferimenti impliciti senza sovraccaricare la
\textit{context window} dell'LLM. La gestione della memoria a lungo termine
e le logiche di \textit{summarization} sono demandate all'Orchestratore
ReAct, circoscrivendo il compito della pipeline al solo
\textit{intent-matching} istantaneo.

% ------------------------------------------------------------------
\subsection*{Fase 3: Analisi Logica della Query (NLP Parsing)}
% ------------------------------------------------------------------
Il testo grezzo e lo storico vengono consegnati al servizio
\path{parse_query_service}, che trasforma il linguaggio naturale in
strutture rigide interrogabili dai sistemi a valle. Il processo opera su
due livelli complementari.\\

\noindent
Il primo livello applica \textbf{spaCy} per l'estrazione deterministica di
elementi ad alta rilevanza: soglie numeriche di budget (es.\ \emph{"sotto i
200 euro"} genera il vincolo \path{max\_price=200}) e indicatori di
condizione (es.\ \emph{"rovinato"} forza \path{condition="usato"}).
Il secondo livello delega a un LLM l'interpretazione delle ambiguità
semantiche, restituendo un oggetto strutturato con marchio, specifiche
tecniche e il campo \path{intent_product_type}: un classificatore binario
che distingue un dispositivo principale (\texttt{Main}) da un accessorio
(\texttt{Accessory}), e che governa il filtraggio nelle fasi successive.\\

\noindent
Fondamentale è l'estrazione sistematica dei \textbf{vincoli di esclusione}
(\path{exclusions}): l'LLM identifica negazioni esplicite o implicite
(es.\ \emph{"senza cerniera"}, \emph{"no ricambi"}) e le struttura in un
insieme di predicati che alimenteranno il rigore dell'Expert Judge nella
Fase~9. In parallelo, il modulo \path{_build_ebay_query} traduce queste
negazioni nella sintassi di esclusione delle API eBay, prependendovi il
segno meno (\texttt{-}), per evitare che termini negativi vengano
interpretati come keyword di inclusione.

% ------------------------------------------------------------------
\subsection*{Fase 4: Ricerca Parallela su API eBay e RAG Locale}
% ------------------------------------------------------------------
Tramite \path{asyncio.gather}, la pipeline lancia simultaneamente due flussi
di recupero su sorgenti eterogenee.\\

\noindent
Il primo interroga le \textbf{Browse API di eBay} tramite \path{search_items},
richiedendo un massimo di 15 inserzioni secondo i criteri costruiti in
Fase 3. Il volume dei risultati è deliberatamente limitato: operare su un sottoinsieme ristretto e pre-selezionato garantisce elaborazioni successive più rapide rispetto a un dataset ampio e rumoroso, ottimizzando contestualmente il numero di chiamate alle API.\\

\noindent
Il secondo attiva il modulo \path{_do_rag}, che implementa un approccio
\textbf{Retrieval-Augmented Generation} sul corpus documentale interno.
La query viene prima espansa semanticamente da \path{expand_query} —
identificando sinonimi e varianti terminologiche — quindi biforcata su due
collezioni Qdrant distinte: \path{product_docs} (schede tecniche) e
\path{seller_docs} (guide operative e problemi ricorrenti segnalati dai
venditori). Questo arricchimento qualitativo contestualizza le inserzioni
eBay con informazioni assenti nelle schede commerciali.

% ------------------------------------------------------------------
\subsection*{Fase 5: Fusione Ibrida dei Risultati}
% ------------------------------------------------------------------
I risultati della ricerca live vengono fusi con il corpus storico del
Phantom Fetch nel modulo \path{_merge_hybrid_results}. Prima della fusione,
ogni inserzione storica è sottoposta a due verifiche: la conformità ai
vincoli di budget correnti e la disponibilità effettiva sull'API eBay
tramite il suo \path{ebay_id}, scartando prodotti venduti o rimossi dalla
prima indicizzazione.\\

\noindent
Il sistema applica poi un meccanismo di rinforzo (\path{_hybrid_bonus}):
i prodotti presenti sia nel corpus storico sia tra i risultati live ricevono
un incremento di \(+0.15\) punti sullo score RRF (Reciprocal Rank Fusion). La logica è diretta: un
prodotto che emerge spontaneamente da due sorgenti indipendenti ha maggiori
probabilità di essere genuinamente rilevante.

% ------------------------------------------------------------------
\subsection*{Fase 6: Elaborazioni Parallele (Profilazione e Trust)}
% ------------------------------------------------------------------
Prima di procedere con interrogazioni intensive di dettaglio, la pipeline esegue due operazioni in background fondamentali per la personalizzazione asincrona del ranking.\\

\noindent
In primo luogo, il modulo \path{update_user_profile} analizza l'intento logico della ricerca per aggiornare dinamicamente il profilo dell'utente (ad esempio per registrare una forte preferenza contingente per un brand o una categoria) all'interno del database relazionale. In secondo luogo, tramite \path{_prefetch_top_sellers_feedback}, viene calcolata ed iniettata in parallelo la reputazione dei venditori in gara, valutando la composizione sentimentale dei loro feedback. L'estrazione tempestiva del \emph{Trust Score} in questa fase è vitale poiché costituisce un segnale chiave per l'architettura vettoriale a valle.

% ------------------------------------------------------------------
\subsection*{Fase 7: Arricchimento Contestuale Profondo}
% ------------------------------------------------------------------
Prima del giudizio semantico definitivo, la pipeline effettua un passaggio
di \textbf{arricchimento totale}. Per ogni identificatore univoco
(\path{ebay_id}) del pool (inserzioni live e RAG) vengono eseguite
chiamate asincrone parallele all'API di dettaglio eBay (\path{getItem}).\\

\noindent
Questo recupera la descrizione integrale e gli aspetti tecnici (Item Specifics) di ciascun prodotto: dati strutturati non
disponibili nei risultati sintetici della Browse API. L'arricchimento è il
prerequisito architetturale della fase successiva: senza le specifiche
ufficiali originali, l'Expert Judge non potrebbe verificare la conformità
ai vincoli di esclusione in modo oggettivo, ricadendo in inferenze
allucinatorie.

% ------------------------------------------------------------------
\subsection*{Fase 8: Vector Reranker e Cross-Encoder}
% ------------------------------------------------------------------
A questo punto della pipeline, il pool di candidati è stato filtrato,
arricchito e fuso da sorgenti eterogenee. Il problema che rimane è
decidere l'ordine: quale prodotto mostrare per primo? Questa fase
risponde a quella domanda attraverso due stadi in cascata, dal più
rapido al più preciso.

\paragraph{Stadio 1 — Base-Score.}
Ogni inserzione viene valutata su un sistema di oltre quindici
\textit{feature} ingegnerizzate manualmente, ciascuna delle quali
cattura una dimensione distinta della rilevanza. Tra le principali
figurano la similarità semantica tra la query e il titolo del prodotto
(\path{semantic_sim}), la similarità lessicale (\path{lexical_sim}),
l'affidabilità del venditore derivata dall'analisi del sentiment dei
feedback (\path{trust_score}), la distanza del prezzo dalla mediana
di mercato (\path{price_z}), il boost derivante dal corpus documentale
RAG (\path{rag_product_boost}) e il punteggio di accessorietà
(\path{accessory_score}). Il risultato è un punteggio numerico
aggregato che riflette la qualità complessiva dell'inserzione rispetto
alla richiesta.

\paragraph{Stadio 2 — Cross-Encoder.}
I dieci candidati con il Base-Score più alto vengono sottoposti a
un'analisi semantica più profonda. A differenza del primo stadio,
che valuta query e prodotto separatamente, un modello Cross-Encoder,
nello specifico, \path{ms-marco-MiniLM-L-6-v2}, pre-addestrato
su task di ranking documentale; li legge \emph{insieme}, in un
unico contesto. Questo consente al modello di valutare non quanto
il prodotto assomigli alla query in astratto, ma quanto la risponda
concretamente. Il costo computazionale più elevato giustifica
l'applicazione selettiva ai soli dieci finalisti.

\paragraph{Fusione degli score.}
I due punteggi vengono combinati in un voto finale pesato.
Al cross-score, che ragiona sull'intento reale della
richiesta leggendo query e prodotto insieme, viene assegnato
il 70\% del peso finale. Al Base-Score, che valuta
criteri oggettivi come prezzo e reputazione del venditore,
il restante 30\%. La scelta riflette una gerarchia
deliberata: la comprensione semantica della richiesta è il
segnale più prezioso, ma i criteri numerici concreti non
vengono ignorati.

% ------------------------------------------------------------------
\subsection*{Fase 9: Expert Judge LLM (Chain-of-Thought Reranking)}
% ------------------------------------------------------------------
Questa fase introduce il validatore finale a tolleranza zero (\path{ltr.py}).
Il modello opera tramite un protocollo \textbf{Chain-of-Thought (CoT)}: per
ogni prodotto nella lista arricchita, esegue un ragionamento critico
esplicito che confronta i vincoli di esclusione estratti in Fase~3 con i
dati tecnici profondi ottenuti in Fase~7. Se la descrizione integrale
specifica la presenza di un attributo escluso dalla richiesta
(es.\ \emph{"full-length zipper"} a fronte di \emph{"senza cerniera"}),
il Giudice assegna uno score di \(-1.0\), eliminando l'inserzione dalla
graduatoria indipendentemente da qualsiasi altro segnale di rilevanza.
I voti del Giudice vengono infine combinati con gli score vettoriali della
Fase 8 e il \path{trust_score} reputazionale per determinare l'ordinamento
definitivo dei risultati.

% ------------------------------------------------------------------
\subsection*{Fase 10: Persistenza Transazionale Preliminare}
% ------------------------------------------------------------------
Per isolare la latenza delle operazioni sul database relazionale dal
percorso critico della classificazione, la memorizzazione permanente viene
avviata prima che la graduatoria finale sia conclusa. Il modulo
\path{_prepare_and_persist_items} opera in un \textit{thread} isolato,
salvando i dati su Postgres. \textbf{Redis} agisce a monte come layer di
deduplicazione: tramite la chiave \path{known_ebay_ids}, intercetta ogni
scrittura e blocca gli identificatori già archiviati, inviando in
\texttt{INSERT} esclusivamente le inserzioni inedite. Questo meccanismo
riduce il carico I/O sul database e annulla il rischio di collisioni.

% ------------------------------------------------------------------
\subsection*{Fase 11: Classificazione Finale e Modificatori Deterministici}
% ------------------------------------------------------------------
Ogni inserzione riceve la sua calibrazione definitiva tramite
\path{_apply_final_ranking}. All'ordinamento prodotto nelle Fasi~8 e~9
vengono sovrapposti modificatori deterministici derivati dal profilo
storico dell'utente. Il \textbf{Brand Affinity Bonus} incrementa di
\(+0.15\) punti le inserzioni di brand verso cui l'utente ha documentato
una predilezione storica. L'\textbf{euristica value-for-money} eroga un
premio di \(+0.10\) punti ai prodotti il cui prezzo si posizioni
tangibilmente al di sotto delle soglie di \path{contextual_budgets},
elevando in classifica le opzioni con il miglior rapporto qualità/prezzo
nel contesto soggettivo dell'acquirente.

% ------------------------------------------------------------------
\subsection*{Fase 12: Metriche IR e Indicizzazione Finale su Qdrant}
% ------------------------------------------------------------------
A conclusione dell'elaborazione, la pipeline calcola le metriche di qualità
del recupero secondo gli standard dell'\textbf{Information Retrieval}:
\textbf{Precision@5}, \textbf{Recall@10} e \textbf{nDCG@10} (Normalized
Discounted Cumulative Gain). Questi indicatori sono resi disponibili per
diagnosi offline sulle performance dell'IR engine.\\

\noindent
Contestualmente, i quindici risultati di vertice vengono inviati a
\path{index_search_items} per l'indicizzazione su \textbf{Qdrant}. Questo
passaggio chiude il ciclo architetturale: i prodotti appena classificati e
validati entrano nel corpus disponibile per i Phantom Fetch delle sessioni
future, migliorando progressivamente la qualità del recupero anticipato su
query semanticamente affini.

% ============================================================
% LTR EXPERIMENT
% ============================================================
\section{Approfondimento di Ricerca: Learning to Rank Personalizzato (LTR)}

Come anticipato nell'evoluzione della pipeline, l'architettura di
classificazione ha attraversato una fase di sviluppo intensivo orientata
all'addestramento di un modello supervisionato nativo. Benché questo approccio
sia stato superato in produzione dall'attuale sistema ibrido, le logiche
simulative e ingegneristiche sviluppate per sopperire alle limitazioni del
dominio costituiscono un contributo architetturale di rilievo che merita
documentazione formale.

\subsection*{Studio Preliminare: Personalità Acquirenti e Selezione degli LLM}

Prima di procedere alla progettazione della simulazione, è stato condotto uno
studio approfondito su due assi complementari: la caratterizzazione delle
personalità acquirenti più diffuse nell'e-commerce online e l'identificazione
dei modelli linguistici che meglio ne incarnano i tratti comportamentali.\\

\noindent
Sul primo asse, la letteratura accademica e i report di settore convergono
nell'identificare cluster psicografici ricorrenti tra i consumatori digitali,
ciascuno con euristiche decisionali, tolleranze al rischio e sensibilità al
prezzo ben distinte. Profili come il \textit{Bargain Hunter}, orientato
esclusivamente alla minimizzazione del costo, o il \textit{Luxury Buyer},
insensibile al prezzo ma esigente sulla reputazione del venditore, non sono
costrutti arbitrari, ma riflettono distribuzioni empiricamente validate nei
clickstream di piattaforme come eBay e Amazon.\\

\noindent
Sul secondo asse, è emerso che i \textit{Foundation Model} differiscono
significativamente nella loro propensione naturale a simulare determinati
profili. Architetture addestrate su corpus prevalentemente anglofoni con
allineamento conservativo (es.\ Llama) tendono a incarnare più fedelmente
profili cauti e orientati alla qualità; modelli con parametri di allineamento
più laschi o addestrati su distribuzioni culturalmente diverse (es.\
DeepSeek, Qwen) mostrano una maggiore varianza nelle valutazioni e una minore
tendenza alla compiacenza organica, rendendoli più adatti a simulare profili
impulsivi o irrazionali. Istanze locali a temperatura elevata (\textbf{Ollama})
sono state invece privilegiate per i profili ad alta entropia decisionale,
come l'\textit{Emotional Bidder}.\\

\noindent
Questo studio ha fornito la base metodologica per la progettazione delle
Persona e per l'assegnazione deliberata dei modelli ai profili nella fase
di simulazione.

\subsection*{Obiettivo e Infrastruttura}

L'obiettivo era addestrare un modello LTR capace di riordinare i risultati
in base al profilo comportamentale dell'utente: comprendere, ad esempio,
che un consumatore orientato al risparmio privilegia sistematicamente le
fasce di prezzo inferiori, indipendentemente dalla pertinenza nominale.
Il problema è stato formulato come ranking supervisionato \textbf{Pointwise},
in cui ogni coppia \emph{(utente, prodotto)} riceveva un'etichetta di
rilevanza da $0$ a $N-1$, con addestramento su \textbf{LightGBM}
(\path{rank:ndcg}). A differenza di un cross-encoder isolato, l'LTR
consentiva il vaglio simultaneo dell'intero pool di candidati e integrava
le \textit{feature utente} direttamente nel peso di predizione.

\subsection*{Il Problema dei Dati e l'Ecosystem Simulation}

L'ostacolo strutturale è stata la mancanza di dataset pubblici di
\textit{clickstream} per eBay. La soluzione è stata la progettazione di un
\textbf{Dataset Sintetico} tramite simulazione algoritmica
(\textit{Ecosystem Simulation}): un layer di 15 agenti LLM incaricati di
impersonare le Persona identificate nello studio preliminare, vincolati a
parametri numerici rigidi (\path{price_sensitivity}, \path{quality_preference}, \path{brand_loyalty}) e a \textit{System Prompt} calibrati per inibire
deviazioni creative e ancorare l'inferenza al profilo assegnato.\\

\noindent
L'assegnazione LLM-Persona non è stata casuale ma motivata dallo studio
precedente: a ciascun profilo è stato abbinato il modello la cui distribuzione
inferenziale naturale risultava più compatibile con le euristiche decisionali
del profilo stesso. Questo ha trasformato i \textit{bias} intrinseci di
ciascuna architettura in uno strumento: il bias
conservativo di \textbf{Llama 3.3} per i profili cauti, la maggiore
varianza inferenziale di \textbf{DeepSeek} e \textbf{Qwen} per i profili
impulsivi, l'alta entropia di istanze \textbf{Ollama} (locali) a temperatura elevata
per i profili irrazionali. I punteggi prodotti venivano infine normalizzati
tramite scale logaritmiche $\log(1+x)/\log(1+p_{99})$ per smorzare le
derive esponenziali di variabili ad alto range come il numero di feedback
per venditore.

\subsection*{Integrazione LightGBM e Conclusioni}
Le estrazioni depurate hanno prodotto un dataset di \textbf{19 feature}
per migliaia di coppie \emph{(query, prodotto)}, memorizzate nel file
\path{training_set.csv}. Prima di addestrare il modello, è stato
necessario verificare che i giudizi sintetici prodotti dalle diverse
Persona fossero sufficientemente coerenti tra loro ma non identici:
una convergenza totale avrebbe indicato che i modelli si stavano
semplicemente copiando, mentre una divergenza eccessiva avrebbe reso
il dataset inaffidabile. Questa coerenza inter-modello è stata misurata
tramite il \textbf{Cohen's Kappa Quadratico Pesato}, uno strumento
statistico standard per valutare l'accordo tra valutatori indipendenti.
Il valore ottenuto, $\kappa \in [0.4, 0.6]$, ha confermato una
eterogeneità positiva: i modelli ragionavano in modo abbastanza simile
da essere utili, ma abbastanza diverso da arricchire il dataset.\\

\noindent
Il dataset è stato quindi utilizzato per addestrare un modello
\textbf{LightGBM} con obiettivo \textit{lambdarank}, ottimizzato
direttamente sulla metrica nDCG — la stessa usata per valutare la
qualità del ranking nella pipeline.\\

\noindent
L'analisi dei risultati ha tuttavia rivelato un problema diverso da quello
atteso. I giudizi sintetici non erano semplicemente distorti verso descrizioni
organiche: erano \textbf{eccessivamente polarizzati}. Le Persona tendevano
ad assegnare valutazioni estreme — punteggi molto alti o molto bassi —
con una deviazione standard dei giudizi prossima a $3$, un valore che
indica scarsa capacità discriminativa nelle fasce intermedie. Il paradosso
è che il \textbf{Cohen's Kappa Quadratico Pesato} risultava comunque
accettabile ($\kappa \in [0.4, 0.6]$): i modelli \emph{concordavano}
su cosa fosse buono o cattivo, ma lo facevano in modo troppo netto,
senza le sfumature graduali che caratterizzano il comportamento reale
di un consumatore. In assenza di un \textit{gold set} organico con cui
ancorare e correggere questa polarizzazione, il dataset si è rivelato
insufficiente ad addestrare un modello di ranking credibile.\\

\noindent
Questo fallimento ha però direttamente orientato la progettazione
dell'\textbf{Expert Judge LLM} descritto in precedenza.